\chapter{Flows}\label{ch:s1:flows}

When investors pull money from a mutual fund, the fund sells what it holds; when many funds lose money together, they sell the same stocks together, and
the prices fall for reasons that have nothing to do with the companies. Coval and Stafford found this \omterm{def:s1:flows:fit}{price pressure} in mutual-fund trades from 1980 to
2003, and found that whoever traded against the forced sellers was paid for it. Lou aggregated the trading that fund flows force into a measure for each
stock and found a significant, temporary price impact, reversed in the following years; since flows are predictable, so is the pressure. This chapter
simulates two hundred funds on the synthetic market, measures the pressure their flows exert and how much of it reverses, and tries the two trades the
literature suggests: getting ahead of the flows, and buying what they forced out. The build is \texttt{firm.flowpress}.

\section{Flow-driven price pressure}

\begin{definition}[Flow-induced trading, price pressure]\label{def:s1:flows:fit}
\emph{Flow-induced trading}\index{flow-induced trading} of a stock over a period is the trading that fund flows force on it: the sum over funds of the share of
the stock each fund owns times the fund's flow as a share of its assets. \emph{Price pressure}\index{price pressure} is the part of a price move caused by
such uninformed demand, expected to reverse once the demand has been absorbed.
\end{definition}

Why should uninformed demand move prices at all? If other investors stood ready to take the other side at a price close to value, it would not. Gabaix and
Koijen argue that the market as a whole is inelastic: institutions operate under mandates that leave little room to absorb flows, and they estimate that
\$1 invested in the stock market raises its value by about \$5. For single stocks, the \omterm{def:s1:flows:fit}{flow-induced trading} of Coval and Stafford and of Lou is the measured
version of the same idea.

The synthetic market's funds (\texttt{firm.flowpress}) are simple. Two hundred funds hold sixty stocks each, picked at random, and replace a quarter of their
holdings every quarter; together they own about 20\% of each stock. Each quarter a fund's flow is 5\% of its assets per standard deviation of its trailing
one-year return relative to other funds, plus 3\% of noise: flows chase performance. A stock's \omterm{def:s1:flows:fit}{flow-induced trading} that quarter (standard deviation 0.5\% of
its shares) pushes its price by 1.5 times itself, spread over the quarter, and the push reverses with a half-life of 126 days. Because the flows chase
returns that include past pressure, pressure breeds flows: the mechanism Lou proposed for fund performance persistence.

\section{Mutual-fund redemptions and fire sales}

Measured in the cross-section, each quarter's log return moves 1.05 per unit of that quarter's \omterm{def:s1:flows:fit}{flow-induced trading}; over the following year, $-0.69$, so that
66\% of the quarter's move reverses within a year (the planted decay alone would take back 75\% of the end-of-quarter push). Without planted pressure the
same slopes are 0.31 and 0.32: funds that did well hold stocks that did well, and the synthetic market's stocks have persistent drifts, so flows point at
drifting stocks whether or not they push them. Any measure of flow pressure has to be read against that confound.

\begin{omfigure}
\begin{tikzpicture}
\begin{axis}[width=11cm,height=5.4cm,xlabel={\small years from the quarter's start},ylabel={\small top minus bottom decile (\%)},
  tick label style={font=\footnotesize},xmin=0,xmax=3,ymin=-4,ymax=4,grid=major,grid style={gray!25},legend pos=north east,
  legend style={font=\scriptsize},table/col sep=comma]
\addplot[omDef,very thick] table[x=year,y=spread]{figdata/strategies-1/09-flows/path.csv};
\addplot[black!50,thick,dashed] table[x=year,y=no_pressure]{figdata/strategies-1/09-flows/path.csv};
\addplot[black!40,dotted] coordinates {(0,0) (3,0)};
\legend{with flow pressure,same funds and flows without pressure}
\end{axis}
\end{tikzpicture}

\omcaption{\omterm{def:s1:earnings:event}{Event study} of \omterm{def:s1:flows:fit}{flow-induced trading} on the synthetic market: cumulative log return of the stocks in the top decile of a quarter's \omterm{def:s1:flows:fit}{flow-induced
trading} minus the bottom decile, from the quarter's start, with the planted pressure and without it. Data: \texttt{s1\_flows.event\_path}.}
\label{fig:s1:flows:path}
\end{omfigure}

The \omterm{def:s1:earnings:event}{event study} (\cref{fig:s1:flows:path}) shows the shape Lou describes. The top decile of a quarter's \omterm{def:s1:flows:fit}{flow-induced trading} beats the bottom by 1.74\% over the
quarter and by 1.91\% at its peak, three days after the quarter ends; the gap is back to zero 261 trading days after the quarter's start, about a year, and ends at $-2.15\%$ after
three years. Without pressure the gap is 0.15\% at the quarter's end and drifts to $-1.41\%$.

\section{ETF and index-fund flows}

Mutual funds trade their flows with some discretion; exchange-traded funds cannot. An ETF's creations and redemptions are baskets of its holdings, and
arbitrageurs trade the basket whenever the ETF's price leaves its value. Ben-David, Franzoni and Moussawi, using changes in index membership as a natural
experiment, found that stocks with higher ETF ownership have significantly higher volatility and more negative autocorrelation (liquidity shocks in the ETF
reach its baskets through the arbitrage) and earn a risk premium of up to 56 basis points a month. The flows of index funds are the most mechanical of all,
and chapter~10 trades the most predictable of them, index rebalancing.

\section{Quarterly holdings as a signal}

\begin{definition}[13F holdings]\label{def:s1:flows:13f}
\emph{13F holdings}\index{13F holdings} are the quarterly equity positions that US institutional investment managers above a size threshold must report to
the SEC on Form 13F, published after a lag; they are the public source for the fund holdings on which flow signals are built.
\end{definition}

\begin{dated}{2026-09}{Form 13F}\label{dat:s1:flows:13f}
Institutional investment managers exercising investment discretion over \$100 million or more in Section 13(f) securities must file Form 13F with the SEC
each quarter, within 45 days after the quarter's end; the SEC publishes the Official List of Section 13(f) securities shortly after each quarter. Mutual
funds also report their full portfolios on their own forms. A holdings-based signal therefore sees positions one and a half to four and a half months old.
\end{dated}

The chapter's two books respect that lag.

\begin{center}
\small
\begin{tabular}{@{}lcc@{}}
\toprule
synthetic funds, years 3 to 10 & with planted pressure & without\\
\midrule
expected flows: Sharpe ratio before, after costs & 0.70, 0.44 & 0.38, 0.12\\
expected flows: return a year after costs; turnover & 1.4\%; 4.0 & ---\\
fire-sale reversal: Sharpe ratio before, after costs & $-0.03$, $-0.35$ & $-0.73$, $-1.06$\\
\bottomrule
\end{tabular}
\end{center}

The \emph{expected-flow} book forecasts next quarter's \omterm{def:s1:flows:fit}{flow-induced trading} at each quarter's start: flows from the funds' trailing returns (with a slope
estimated on past quarters only) times holdings from the previous quarter-end, already filed. It buys the top decile and sells the bottom for the quarter,
and nets a Sharpe ratio of 0.44. The \emph{fire-sale reversal} book waits until a quarter's trading is known from the next filing, 45 days after the quarter,
buys the stocks the flows sold hardest and sells those they bought, and holds for a year. It loses: by the time the trading is known, a good part of the
reversal has happened, and what remains is fought by the drifts the flows were chasing. The version without pressure loses more, which is the pressure's
contribution. The reversal is real in the \omterm{def:s1:earnings:event}{event study} and hard to capture after the lag.

A last variant shows the confound at its worst. Give each fund a style tilt of its own, and flows become style bets: funds tilted to a rewarded style do
well, attract inflows, and buy that style. The top-minus-bottom decile of \omterm{def:s1:flows:fit}{flow-induced trading} then gains 4.9\% in the quarter and 69.0\% over three years,
none of it reversing, because it is the style premium, not the pressure. A flow signal must be checked against the factors it might be proxying.

\section{Strategy files}

\begin{strategyfile}{Fire-sale reversal}\label{strat:s1:flows:firesale}
\sfield{Who pays you, and why} Funds forced to sell by redemptions, who accept a price below value for immediacy.
\sfield{Instruments and venues} Stocks held in common by funds with large outflows.
\sfield{Signal} Flow-induced selling, from fund flows and holdings, as soon as it can be known.
\sfield{Sizing and execution} Buy the most-sold names and hold while the pressure reverses; hedge factors.
\sfield{Costs} Moderate turnover; the names are under selling pressure, so liquidity is available to a buyer.
\sfield{How it dies} Reporting lags that let the reversal happen before the signal arrives; confounds with momentum.
\sfield{Horizon, capacity, infrastructure} Months to a year; fund flow and holdings data.
\sfield{Backtest honestly} Holdings and flows as of their publication; control for the factors and momentum.
\sfield{Sources} Coval and Stafford (2007): significant returns to trading against constrained funds, 1980--2003; this chapter's simulation.
\end{strategyfile}

\begin{strategyfile}{Flow-induced pressure}\label{strat:s1:flows:pressure}
\sfield{Who pays you, and why} Fund investors whose flows chase past performance, and the funds that must trade them.
\sfield{Instruments and venues} Stocks held by funds likely to receive or lose money.
\sfield{Signal} Expected \omterm{def:s1:flows:fit}{flow-induced trading}: flows forecast from trailing fund returns times reported holdings.
\sfield{Sizing and execution} Long expected buying, short expected selling, for the quarter.
\sfield{Costs} Quarterly turnover.
\sfield{How it dies} Flows that stop chasing performance; front-runners crowding the same forecast.
\sfield{Horizon, capacity, infrastructure} A quarter; holdings, flows and fund returns.
\sfield{Backtest honestly} Holdings lagged to their filing; the flow model estimated on past data only.
\sfield{Sources} Lou (2012): the expected part forecasts next year's returns, reversed later; Coval and Stafford (2007) on front-running forced trades.
\end{strategyfile}

\begin{strategyfile}{ETF creation-redemption flow}\label{strat:s1:flows:etf}
\sfield{Who pays you, and why} ETF liquidity traders whose trades reach the basket through arbitrage.
\sfield{Instruments and venues} ETFs and their constituents.
\sfield{Signal} Daily creations and redemptions, or the ETF's premium to its value.
\sfield{Sizing and execution} Provide liquidity in the basket against the flow; unwind as it reverts.
\sfield{Costs} Basket trading; short horizons.
\sfield{How it dies} Arbitrageurs faster than the strategy; flows that carry information.
\sfield{Horizon, capacity, infrastructure} Days; ETF flow and holdings files.
\sfield{Backtest honestly} Creation data with its publication time; basket costs.
\sfield{Sources} Ben-David, Franzoni and Moussawi (2018): higher volatility and more negative autocorrelation with ETF ownership.
\end{strategyfile}

\begin{strategyfile}{Holdings-overlap signal}\label{strat:s1:flows:overlap}
\sfield{Who pays you, and why} As for fire sales: funds holding the same stocks sell them together.
\sfield{Instruments and venues} Stocks with concentrated common ownership.
\sfield{Signal} The overlap of a stock's holders with funds under stress (Book~7, chapter~28's overlap measure).
\sfield{Sizing and execution} As a risk flag, or as a conditional reversal trade after a stress.
\sfield{Costs} Low as an overlay.
\sfield{How it dies} Holdings too stale to show current overlap.
\sfield{Horizon, capacity, infrastructure} Quarters; holdings data.
\sfield{Backtest honestly} Overlap from filed holdings only.
\sfield{Sources} Coval and Stafford (2007) on \omterm{def:s1:flows:fit}{price pressure} in securities held in common; no performance figure verified.
\end{strategyfile}

\section{Tutorial: forced sellers}\label{tut:s1:flows:tut}

\begin{tutorial}
\textbf{Goal.} Simulate funds with performance-chasing flows, measure the pressure and its reversal, and trade the expected flows and the fire sales with
the filing lag. \textbf{End state:} the table and \cref{fig:s1:flows:path}.
\begin{tutsteps}
\item \textbf{The funds}: each quarter, flows from trailing returns, \omterm{def:s1:flows:fit}{flow-induced trading} from holdings, and the price push spread over the quarter.
\omcode{code/firm/flowpress/firm_flowpress.py}{90}{105}{Flows, flow-induced trading and the pressure.}\label{lst:s1:flows:sim}
\item \textbf{The expected-flow book}: a flow model estimated on past quarters, and holdings as filed.
\omcode{code/strategies-1/09-flows/python/s1_flows.py}{100}{108}{Forecasting next quarter's flow-induced trading.}\label{lst:s1:flows:book}
\item \textbf{Run} \texttt{response}, \texttt{event\_path}, \texttt{books} (with and without pressure) and \texttt{fig\_flows.py}.
\end{tutsteps}
\textbf{What to change next.} Let fund sizes follow their flows and watch pressure compound; add ETFs that trade their baskets mechanically; shorten the
filing lag and see how much of the reversal a faster signal captures.
\end{tutorial}

\section{Build: fund flows and pressure}\label{bld:s1:flows:flowpress}

\begin{build}
\bfield{Purpose} Simulated fund holdings and flows, \omterm{def:s1:flows:fit}{flow-induced trading}, and the \omterm{def:s1:flows:fit}{price pressure} it causes, to study flow strategies with a known truth.
\bfield{Interface} \texttt{FlowConfig}, \texttt{fit(own, flow)}, \texttt{flow\_rule(trailing, beta, noise, rng)}, \texttt{simulate\_flows(R, listed, style,
cfg)}.
\bfield{Rules} Flows from information available at the quarter's start; the pressure added to returns without touching the rest of the market; holdings
recorded quarter by quarter.
\bfield{Acceptance tests} \texttt{code/firm/flowpress/tests/}: \omterm{def:s1:flows:fit}{flow-induced trading} by hand; flows ranked by trailing returns; the pressure follows the
quarter's trading, and the added log returns sum to the pressure level.
\bfield{Stretch} Fund sizes that follow flows; partial scaling of flows into trades (funds using cash); ETFs.
\end{build}

\begin{omsources}
\item J.~Coval and E.~Stafford, ``Asset fire sales (and purchases) in equity markets'', \emph{Journal of Financial Economics} 86(2), 2007.
\item D.~Lou, ``A flow-based explanation for return predictability'', \emph{Review of Financial Studies} 25(12), 2012.
\item I.~Ben-David, F.~Franzoni and R.~Moussawi, ``Do ETFs increase volatility?'', \emph{Journal of Finance} 73(6), 2018.
\item X.~Gabaix and R.~S.~J.~Koijen, ``In search of the origins of financial fluctuations: the inelastic markets hypothesis'', NBER working paper 28967,
2021.
\item US SEC, Frequently Asked Questions about Form 13F.
\end{omsources}

\section{Exercises}

\begin{exercise}[$\star$]\label{exo:s1:flows:1}
Funds own 20\% of a stock; they receive inflows of 5\% of their assets. What is the stock's \omterm{def:s1:flows:fit}{flow-induced trading}, and what does an impact of 1.5 make of it?
\end{exercise}

\begin{exercise}[$\star$]\label{exo:s1:flows:2}
A push decays with a half-life of 126 days. How much of it remains after 126 days, and after a year of 252?
\end{exercise}

\begin{exercise}[$\star$]\label{exo:s1:flows:3}
A 13F filing for the quarter ending March 31 is due within 45 days. How old are the holdings when the filing arrives, and when the next one does?
\end{exercise}

\begin{exercise}[$\star\star$]\label{exo:s1:flows:4}
Why do \omterm{def:s1:flows:fit}{flow-induced trading} and past returns point at the same stocks on the synthetic market even without any pressure?
\end{exercise}

\begin{exercise}[$\star\star$]\label{exo:s1:flows:5}
Explain how performance-chasing flows can create fund performance persistence.
\end{exercise}

\begin{exercise}[$\star\star$]\label{exo:s1:flows:6}
Why does the fire-sale reversal book lose even though the \omterm{def:s1:earnings:event}{event study} shows a reversal?
\end{exercise}

\begin{exercise}[$\star\star\star$]\label{exo:s1:flows:7}
\emph{Coding.} Run \texttt{event\_path(1.5, 756, 2.0)}. Explain why tilted funds turn the flow signal into a style bet, and how you would test a real flow
signal for it.
\end{exercise}

\begin{exercise}[$\star\star\star$]\label{exo:s1:flows:8}
\emph{Find the flaw.} ``We compute \omterm{def:s1:flows:fit}{flow-induced trading} from quarter-end holdings and trade on the quarter-end date; the backtest earns 8\% a year.''
\end{exercise}

\section{Problem: Forced Sellers}

\begin{problem}[{Weekend problem --- pressure, reversal and the lag}]\label{pb:s1:flows:1}
The chapter's simulated funds on \texttt{firm.synthmkt}, and the public record.

\medskip\noindent\textbf{Part I --- The idea.}
\begin{enumerate}
\item Define \omterm{def:s1:flows:fit}{flow-induced trading} and \omterm{def:s1:flows:fit}{price pressure}.
\item What did Coval and Stafford find?
\item What did Lou find, and what does it explain?
\item What is the inelastic markets hypothesis?
\end{enumerate}

\medskip\noindent\textbf{Part II --- The simulation.}
\begin{enumerate}[resume]
\item Describe the funds, their flows and the planted pressure.
\item How does the pressure feed back into flows?
\item Give the cross-sectional responses with and without pressure.
\item What share of the quarter's move reverses within a year?
\end{enumerate}

\medskip\noindent\textbf{Part III --- The trades.}
\begin{enumerate}[resume]
\item Describe the expected-flow book and its result.
\item Describe the fire-sale reversal book and why it loses.
\item What do ETFs change?
\item What is Form 13F, and what lag does it impose?
\end{enumerate}

\medskip\noindent\textbf{Part IV --- The verdict.}
\begin{enumerate}[resume]
\item State the \emph{named result}: the \omterm{def:s1:flows:fit}{price pressure} of \omterm{def:s1:flows:fit}{flow-induced trading} and the share that reverses within a year.
\item Read the \omterm{def:s1:earnings:event}{event study}.
\item What does the style-tilted variant show?
\item Where would you get faster flow data?
\item How would you control a flow signal for momentum?
\item Which strategy file would you run?
\item What would make the fire-sale trade work?
\item In one sentence: what does a flow strategy sell?
\end{enumerate}
\end{problem}

\section{Interview questions}

\begin{interviewq}[$\star$ \iqroles{researcher}]\label{iq:s1:flows:1}
Why would mutual-fund flows move stock prices?
\end{interviewq}

\begin{interviewq}[$\star\star$ \iqroles{researcher, developer}]\label{iq:s1:flows:2}
How would you compute \omterm{def:s1:flows:fit}{flow-induced trading} for every US stock from public data, and with what lag?
\end{interviewq}

\begin{interviewq}[$\star\star$ \iqroles{researcher}]\label{iq:s1:flows:3}
Your flow signal predicts returns. How do you show it is not momentum in disguise?
\end{interviewq}

\begin{interviewq}[$\star\star$ \iqroles{trader}]\label{iq:s1:flows:4}
A large fund is known to be liquidating. How would you trade it, and what would you worry about?
\end{interviewq}

\begin{interviewq}[$\star\star$ \iqroles{risk}]\label{iq:s1:flows:5}
What does common ownership by funds under stress mean for a stock's risk?
\end{interviewq}

\begin{interviewq}[$\star\star\star$ \iqroles{researcher}]\label{iq:s1:flows:6}
A stock is pushed by $\lambda F$ over a quarter and the push decays exponentially with half-life $h$ days afterwards. Write the expected return of a strategy
that buys at the quarter's end plus a lag of $L$ days and holds for $H$ days, and find when it is positive.
\end{interviewq}
